\section{A native nonhyperfinite family with a half-rate bound}
\label{sec:ne}

The directed recovery formulation does not require a decomposition into
small components.  We demonstrate this with one fixed native computation.
Its full cone incidence graphs contain robust expander cores, yet every
storage code on every induced coordinate set has logarithmic cardinality
at most half its coordinate count plus $o(L)$.  This gives rate at most
one half plus $o(1)$ whenever the retained set has size $\Omega(L)$.
The computation is an invertible shear,
not a transposition algorithm.  The point of the example is the simultaneous
graph-theoretic and coding conclusion for its actual execution.

Throughout this section let
\[
 d=2m,\qquad b=\max\{1,\lceil\log_2d\rceil\},\qquad
 W=db,\qquad L=dW=d^2b.
\]
A raw coordinate is $p=(i,j,s)\in V=[d]\times[d]\times[b]$, where
$[u]=\{0,\ldots,u-1\}$.  Its physical address is $(id+j)b+s$, and its
transpose demand is the physical output address of
$\Pi(p)=(j,i,s)$.  The paired incidence graph has distinct ports $L_p,R_p$;
an edge $L_pR_q$ means that raw input $p$ occurs in the full backward cone
of the append at $\Pi(q)$.  We write this graph as $E_\Gamma$ when cones
are stopped by a whole-step genie $\Gamma$.  Deleting a coordinate always
deletes both of its ports.  In particular it is different from deleting
chosen dependencies.

\subsection{Two imported interfaces and a paid executor}

We use the following expander input from the native shear construction of
\cite[section on expander shears and local recovery]{BrendeLocalCoding}.
The graph construction is imported; the schedule refinements and the
subsequent cone analysis will be given here.

\begin{lemma}[Ordered expander input]\label{lem:ne-expander-import}
\cited{The expander-shear construction of~\cite{BrendeLocalCoding}. The uniform list-generation refinement is justified below by a written argument; this full interface is not claimed to be Lean-verified.}
There are fixed positive constants $J,\Delta,\eta$, with $J,\Delta$
integers and $\eta\le1$, and simple bipartite graphs $H_m$ on
\[
 A_0,\ldots,A_{m-1},\quad B_0,\ldots,B_{m-1},
\]
such that their maximum degree is at most $\Delta$, their edges partition
into $J$ partial increasing matchings, and
\[
 |N_{H_m}(S)\setminus S|\ge\eta|S|
 \quad\text{whenever }|S|\le m.
\]
The matching endpoint lists can be generated uniformly in
$O(m\operatorname{polylog}m)$ time using a fixed number of tapes and
polylogarithmic working space besides the monotonically emitted lists.
All constants and rational coefficients are fixed before $m$ varies.
\end{lemma}

Here a partial increasing matching has its two endpoint sequences strictly
increasing.  The computational part of the imported interface can be seen
directly from the rational overlap construction.  Fix its finite family
of increasing rational maps on interval domains.  Their derivatives have
fixed positive lower and upper bounds.  An edge records a positive-length
overlap of the image of one mesh cell with a second mesh cell; exact
rational comparisons decide boundary coincidences.  There are constantly
many candidate neighbors in either direction.  Split candidates by map
name and by the residues of the two endpoint indices modulo a fixed
integer larger than these consecutive-neighbor bounds.  Each resulting
class is an increasing matching.  Symmetrization uses further fixed
classes, and retaining a duplicate only in its first class preserves the
simple graph.  Evaluation, division, comparison, and duplicate tests use
polynomially many bit operations and cells in $\log m$ per candidate.
Enumerating sources and filtering by class emits each page in increasing
order.  Thus this use of the expander input requires neither a table
depending on $m$ nor a unit-cost operation on a growing word.

The depth interface is a machine-independent consequence of the previous
paper.  We state precisely the portion needed below.

\begin{lemma}[Cheap whole-step depth for home-based walks]
\cited{Theorem B and the proof of Corollary B of~\cite{BrendeCurved}. Both the home-based cut theorem and the whole-step conversion used here are written proofs, not Lean-verified results.}
\label{lem:ne-depth-import}
Fix an oblivious native machine with $k$ work tapes, whole-step value width
at most $\beta$ bits, and running time at most $c(L+1)$.  Suppose every
work-head walk is home-based: every greedy weakly monotone leg begins or
ends within distance one of its fixed home cell.  For every fixed
$\varepsilon>0$, and all sufficiently large admissible $L$, there is a
whole-step genie satisfying $\beta|\Gamma|\le\varepsilon L$ whose full
depth, in the sense of \cref{def:nb-cones}, is at most a constant $D$.
The constant may
depend on the machine and $\varepsilon$, but not on $L$.
\end{lemma}

This is the cut theorem and long-edge conversion of
\cite[Theorem B and the proof of Corollary B]{BrendeCurved}, before any
assumption of transposition correctness.  A cut edge $t\to s$ is blocked
by revealing $s$; a periodic set of additional revealed steps bounds the
runs of state edges and pauses between long hops.  The imported statement
therefore controls all edges of the whole-step DAG, not just long hops.
We will use it only for existence.  Our conclusions about a specified
depth-$D$ genie quantify over every such genie.

\begin{proposition}[A fixed native shear executor]\label{prop:ne-native-shear}
\notformalized
There is a fixed deterministic oblivious native machine $M_0$ which, at
all sufficiently large even canonical dimensions, computes
\[
 \Phi(A,B)=(A,B+L_mA)
\]
in $O(L)$ transitions, where $A,B$ are the first and last $m$ raw rows and
$L_m$ is the binary adjacency matrix of $H_m$.  It uses read-only raw
input, blank initial work tapes, and append-only raw output, and all work
heads are home-based.  There is also such a machine $M_1$ which appends
all $A$ rows only after completing the shear pages, followed by all $B$
rows.  Both machines have fixed alphabet, state set, and tape count.
\end{proposition}

\begin{informalproof}
We give the execution details because later arguments use its visits,
including visits at which the scanned symbols are semantically irrelevant.
Two separate work tapes will hold banks of length $n=mW$.  They begin
blank; each receives a persistent origin flag, row-start flags during
loading, and an end marker when its loading finishes.  A unary row clock
has length $W$ and its own origin and end marker.  All flags are tracks
of a fixed finite alphabet.  Both data heads remain at their origins
throughout setup, so setup does not visit or initialize the full banks.

First count the raw input length with a little-endian binary counter whose
head returns to its origin after each carry.  The sum of carry lengths is
$O(L)$, including return traversals.  The raw input is then rewound.
Canonical shape recovery uses polynomially many bit operations in
$\log L$: try the possible entry widths, compute an integer square root,
and check the defining equality.  Invalid lengths, odd dimensions, and a
fixed finite set of small cases use an identity fallback.  Only the even
canonical subfamily is used in this section.

All polylogarithmic-space arithmetic can be made physically home-based
at polynomially larger polylogarithmic cost.  To make this explicit,
bootstrap a unary scratch ruler of length $S=(\lceil\log_2(L+2)\rceil+c)^h$
by repeated complete sweeps.  Here the fixed exponent $h$ exceeds the
space exponent of the routines to be simulated.  A fixed number of tracks
in this window represent virtual tapes, symbols, and virtual head flags.
A complete outward-and-return scan reads the flagged symbols; further
complete scans implement a virtual transition and relocate flags.  Right
moves can be flagged on an outward scan and left moves on a return scan,
using old and new flag tracks.  Each virtual transition costs $O(S)$,
and every physical turn occurs at the origin or at the end of the ruler.
The length counter itself is implemented by the direct carry routine,
so this simulation does not multiply its linear cost by a polylogarithm.

Generate the $J$ endpoint lists of \cref{lem:ne-expander-import} as
monotone output streams.  Copy an individual record to the origin buffer
before parsing it.  The long list heads make monotone traversals and full
rewinds, rather than reversing locally for arithmetic.  Their total length
is $O(m\log m)$ and the compiler still costs
$O(m\operatorname{polylog}L)$.  Construct the unary row clock by reading
the binary digits of $W$ from most to least significant.  At each digit,
double the current unary string onto an alternate tape, append the digit,
erase and rewind the old string, and rewind the new one.  The geometric
sum of the string lengths is $O(W)$.  All these work heads are home-based.

Load the $A$ bank and then the $B$ bank in raw row order.  The row clock
identifies the first bit of each row, where the load writes the row-start
flag; completion of a bank writes its end marker at $n$.  Thus these
flags require no preliminary full-bank sweep.  An active bank head and
the row-clock head advance one cell per bit.  At a row boundary
the bank head waits while the clock returns to its origin.  The machine
$M_0$ also appends each $A$ bit at its load transition; $M_1$ omits these
early appends.  After loading, rewind both banks completely.

For each increasing matching page, read its next endpoint pair $(i,j)$,
skip whole rows forward until the two indicated rows are reached, and
scan the pair in alignment.  In a single transition at offset $s$, read
$A_i[s]$ and $B_j[s]$, write their XOR to $B_j[s]$, and advance both heads
and the row clock.  A skip likewise uses one transition per bit and
preserves its payload.  Clock rewinds and controller operations take place
with the bank heads parked at row starts or at the bank end.  Increasing
endpoints make a single forward traversal sufficient.  After the last
matched pair, scan the unused suffixes to the end and rewind both banks
in lockstep, one cell per transition.  Empty pages use the same complete
traversal convention.  Each page costs $O(L)$.

All work cells scanned on every transition are rewritten, including cells
under parked heads.  Preserving a symbol is a write, and no part of the
construction assumes otherwise.  Induction over the pages shows that the
$B$ bank contains $B+L_mA$: every simple edge is executed exactly once,
so duplicate XORs do not cancel a desired contribution.  The $A$ bank
remains unchanged.  Finally $M_0$ appends the $B$ bank.  For $M_1$, after
the last full rewind, scan and append $A$ and then $B$.  This has the
specified delayed output order.

The complete cost is
\[
 O\bigl(L+W+JL+m\operatorname{polylog}L+\operatorname{polylog}L\bigr)
 =O(L),
\]
since $L=4m^2b$ and all logarithmic powers are fixed.  Public parameters,
flags, and endpoint records determine every move and append time; payload
values determine only written or appended bits.  The schedule is therefore
oblivious.  Bank and metadata heads make complete outward traversals and
returns, possibly pausing along a traversal.  Clock and compiled controller
heads also make origin excursions.  These are home-based physical walks.
The number of pauses need not be bounded, and we make no assertion of a
bounded number of affine pieces within a leg.
\end{informalproof}

\subsection{Literal expander paths and robust separators}

The native DAG has a state edge $t\to t-1$ and, for each work tape, an
edge from $t$ to the previous visit to its scanned cell, if one exists.
Its in- and out-degrees are at most $k+1$ after parallel edges are merged.
Raw input bits attach to their read steps; the read-only input tape does
not contribute previous-visit edges.  A genie-free path may start at a
revealed vertex but has no subsequent revealed vertex.  When proving
survival below, we conservatively discard a path meeting the genie even
at its first vertex.

\begin{lemma}[A literal core with bounded congestion]\label{lem:ne-paths}
\notformalized
For either machine in \cref{prop:ne-native-shear}, every edge
$A_iB_j$ of $H_m$ and every lane $s\in[W]$ has a literal backward DAG
path from the append of $B_j[s]$ to a read of raw input $A_i[s]$.
The chosen paths have maximum vertex congestion at most $C_0=2\Delta$,
including endpoints.  For lanes $1\le s<W$, their lengths are bounded by
a constant depending only on $J$.  No length bound is needed for lane zero.
\end{lemma}

\begin{informalproof}
Let $t_{\mathrm{in}}$ be the load of $A_i[s]$, $t_{\mathrm{xor}}$ the update for the
chosen edge, and $t_{\mathrm{out}}$ the append of $B_j[s]$.  List every visit
to the physical $A_i[s]$ cell from $t_{\mathrm{in}}$ through $t_{\mathrm{xor}}$;
then list every visit to the $B_j[s]$ cell from $t_{\mathrm{xor}}$ through
$t_{\mathrm{out}}$.  Consecutive visits in each list are actual previous-visit
edges when read backward, and the lists share the XOR vertex.  Their
concatenation gives the claimed path to an actual raw read.  It includes
unchanged writes and controller waits.

For $s\ne0$, no bank head parks at this cell.  The $A$ cell has at most
$2J+2$ relevant visits and the $B$ cell at most $2J+3$; a bound of
$4J+8$ on path length is more than sufficient.  At lane zero, all parked
visits are included and can make the path long.  At any time a scanned
$A$ cell belongs to paths for at most $\Delta$ incident expander edges,
and the same holds for the scanned $B$ cell.  Their union has size at
most $2\Delta$.  The extra delayed $A$ scan in $M_1$ can add parked
visits on the $B$ origin, but does not change this congestion estimate.
\end{informalproof}

Thus $E_\Gamma$ contains $W$ vertex-disjoint copies of $H_m$ with at most
$C_0|\Gamma|$ missing core edges.  There are $2mW=L$ core ports, a subset
of the $2L$ incidence ports.  The claim concerns containment in the full
syntactic cone graph; unselected state and cell paths may supply further
edges.

\begin{proposition}[A robust separator estimate]\label{prop:ne-separator}
\notformalized
Let $U$ be any deleted coordinate set.  If a further port set $Z$ makes
every component of $E_\Gamma-U-Z$ have size at most $K\le m$, then
\[
 |Z|\ \ge\ \frac{\eta L-2C_0|\Gamma|}{\Delta+\eta}-2|U|.
\]
This holds for every genie, without a depth hypothesis.
\end{proposition}

\begin{informalproof}
Discard the at most $f=C_0|\Gamma|$ core edges whose paths meet $\Gamma$.
At most $z\le2|U|+|Z|$ core ports have been deleted.  Intersect each
remaining component with each lane core.  These intersections have size
at most $m$, so expansion applies to each.  Summing their edge-boundary
lower bounds counts at least $\eta(L-z)$ crossing edges with multiplicity.
An edge joining two surviving intersections must be one of the $f$ lost
core edges, and is counted at most twice.  Edges meeting the deleted core
ports contribute at most $\Delta z$.  Hence
\[
 \eta(L-z)\le2f+\Delta z.
\]
Rearrange and use $z\le2|U|+|Z|$.  Additional cone edges can merge
components, but cannot invalidate this lower bound.
\end{informalproof}

\subsection{The full syntactic cones retain a lane constraint}

The preceding core alone says nothing about code rate.  For that purpose
we must control every cone edge, including those introduced by parked
heads and state transitions.  Fix either $M_0$ or $M_1$, let $k$ be its
number of work tapes, and write
\[
 B_D=\sum_{h=0}^D(k+1)^h.
\]

\begin{lemma}[A sparse interface for full cones]\label{lem:ne-lane}
\notformalized
There are fixed constants $C_s,h_s$ and a set $X$ of actual time vertices,
depending only on the length-specific schedule, such that
\[
 |X|\le S_d:=C_s\bigl[W+d(1+\log_2(L+2))^{h_s}\bigr].
\]
For every full-depth-$D$ genie, every output outside the $D$-step forward
neighborhood of $X$, and every non-row-start raw input in its cone, the
two physical within-row bit offsets have cyclic distance at most $D$
in $\mathbb Z/W\mathbb Z$.
\end{lemma}

\begin{informalproof}
The entire length count, shape calculation, metadata generation, and clock
construction form the setup.  This can take $\Theta(L)$ transitions, so
it cannot be discarded as a small set of times.  Its distinct work cells,
however, number only $O(W+d\operatorname{polylog}L)$.  Mark the first
post-setup state vertex and the first post-setup visit to every cell touched
during setup.  Any backward state or cell edge entering setup begins at
one of these marked vertices.

Post-setup fast phases are a forward row load, skip, or XOR scan; a clock
return; a simultaneous whole-bank rewind; and a final forward output
scan.  Mark all other controller, parser, arithmetic, and metadata ticks.
Their total is $O(d\operatorname{polylog}L)$.  Mark also the first and last
ticks of each fast phase, moving bank visits to offset zero modulo $W$,
and moving clock visits to its endpoints.  There are $O(d)$ such marks,
with the fixed number of pages absorbed in the constant.  A stationary
head at offset zero is not marked at every tick.

Every unmarked post-setup vertex has a lane: the position of its moving
clock modulo $W$, if the clock moves, and otherwise that of its moving
bank.  Simultaneously moving banks have the same lane.  A state edge
between unmarked vertices changes lane by at most one.  For a cell edge,
a stationary head has its immediately preceding time as predecessor once
the first fast tick has been excluded.  A moving bank head at a nonzero
offset has no parked visits at that cell; its preceding post-setup moving
visit has the same lane.  The same is true of a moving interior clock
cell, in a forward row scan or return.  These cases cover all work tapes,
including tapes whose fields the current operation ignores.  Consequently
every edge between good vertices changes lane by at most one cyclically.

The $D$-step forward neighborhood of $X$ contains at most $B_D|X|$ time
vertices, since forward degree is at most $k+1$.  An output outside it
has no marked vertex in a depth-$D$ free cone.  Its cone also cannot enter
setup, by the first-touch marks.  During loading, a non-row-start raw bit
is read only at a load vertex with that same lane.  During row-boundary
pauses the raw head holds the next row-start bit or is outside the payload;
after loading it is outside the payload.  Excluding the $d$ raw row starts
therefore accounts for every remaining input attachment.  Summing lane
changes along a free path proves the result.  The output itself is included
in the forward-neighborhood exclusion, which covers the convention allowing
a free path to start at a revealed vertex.
\end{informalproof}

The following elementary code argument will be useful for both examples.

\begin{lemma}[An order of port colors]\label{lem:ne-color-injection}
\formal{ColorOrder.lean}{color_order_half_bound}
Let a paired incidence graph on $N$ coordinates admit colors
$\alpha(v)$ at $L_v$ and $\zeta(v)$ at $R_v$ such that every edge is
monochromatic and $\alpha(v)\ne\zeta(v)$ for every $v$.  Every storage
code over an alphabet of cardinality $a$ has size at most
$a^{\lfloor N/2\rfloor}$.
\end{lemma}

\begin{outline}
Order the finitely many used colors.  Record the coordinates $v$ for which
$\alpha(v)$ precedes $\zeta(v)$.  Process left-color classes in increasing
order.  If a coordinate was not recorded, then its right color precedes
its left color.  Every coordinate in its recovery neighborhood has that
earlier left color, by monochromaticity, and so has already been recovered.
This proves that projection onto the recorded coordinates is injective on
the code; no extension of an arbitrary partial assignment is needed.
Compare the color order with its reverse. Each coordinate is recorded in
exactly one of them, so one record has at most $\lfloor N/2\rfloor$
coordinates, which proves the claim.
\end{outline}

\begin{theorem}[Half rate on every induced full-cone relation]
\notformalized
\label{thm:ne-full-cone-half}
Fix $D\ge0$.  For either native machine in
\cref{prop:ne-native-shear}, for every whole-step genie of full
depth at most $D$, every $Q\subseteq V$, and every storage code
$\mathcal C\subseteq\mathcal A^Q$ for the induced actual cone relation,
where $a=|\mathcal A|\ge2$ and $N=|Q|$, there is $u\le\min(N,e_D(d))$
such that
\[
 |\mathcal C|\le a^{\lfloor(N+u)/2\rfloor}.
\]
One permissible error bound, uniform in the genie, alphabet, and $Q$, is
\[
 \begin{split}
 g&=\lceil\sqrt d\rceil,\qquad r=\lceil d/g\rceil,\\
 e_D(d)&=B_DS_d+d+2d(2D+1)r+bdg=o_D(L).
 \end{split}
\]
There is no bound on the cost of the genie in this theorem.
\end{theorem}

\begin{informalproof}
Partition the column indices into $r$ consecutive buckets of at most $g$
columns each, and view their bit lanes as consecutive arcs of
$\mathbb Z/W\mathbb Z$.  Guard all lanes within cyclic distance $D$ of
a bucket boundary.  There are at most $(2D+1)r$ guarded bit lanes.
Delete from $V$ the following coordinate sets: demand labels of the bad
physical outputs in \cref{lem:ne-lane}; the $d$ raw row starts; coordinates
whose input or transpose-output lane is guarded; and coordinates $(i,j,s)$
whose row and column belong to the same bucket.  The first set costs at
most $B_DS_d$, the third at most $2d(2D+1)r$, and the fourth at most $bdg$.
The total deletion is at most $e_D(d)$.

On the surviving ports give $L_{(i,j,s)}$ the bucket of $j$, and give
$R_{(i,j,s)}$ the bucket of $i$.  A surviving incidence edge has equal
colors by \cref{lem:ne-lane}: a path of lane length at most $D$ cannot
cross a bucket boundary while both endpoint lanes avoid its guard.
The own-port colors differ by the fourth deletion.  The graph therefore
satisfies \cref{lem:ne-color-injection}.

Intersect the deleted set with $Q$, obtaining $u$ coordinates.  Condition
their values in the original code.  Each nonempty fibre is a storage code
on the induced surviving relation: the removed recovery inputs have fixed
values in that fibre.  \cref{lem:ne-color-injection} bounds its size by
$a^{\lfloor(N-u)/2\rfloor}$.  There are at most $a^u$ fibres, giving
$a^{u+\lfloor(N-u)/2\rfloor}=a^{\lfloor(N+u)/2\rfloor}$.
For fixed $D$, each displayed error divided by $L=d^2b$ tends to zero.
All exclusions were chosen from the schedule and the full relation; the
argument is simultaneous for all alphabets, recovery functions, and codes.
\end{informalproof}

\begin{corollary}[Nonhyperfinite bounded-degree cones with half rate]
\notformalized
\label{cor:ne-combined}
For one fixed native home-based linear-time machine there are constants
$\varepsilon,\alpha>0$ and fixed integers $D,M$ with the following
property.  At every sufficiently large even canonical dimension the class
of genies satisfying $\beta|\Gamma|\le\varepsilon L$ and full depth at
most $D$ is nonempty.  For every such genie, one can delete coordinates
to obtain an induced paired graph on $N=\Theta(L)$ coordinates, of maximum
degree at most $M$, with no own-port edge, such that:
\begin{enumerate}
\item for every fixed $K$ and all sufficiently large $d$, every port set
whose deletion leaves components of size at most $K$ has size at least
$\alpha L/2$;
\item every storage code on every further induced coordinate set has
logarithmic cardinality at most half its coordinate count plus $o_D(L)$,
uniformly over all finite alphabets.
\end{enumerate}
\end{corollary}

\begin{informalproof}
Choose $\gamma>0$ with $\gamma\le\eta/(8C_0)$, and then
$0<\varepsilon\le\min(1,\beta\gamma)$.  Set
$\alpha=\eta/[2(\Delta+\eta)]$.  \cref{prop:ne-separator}
gives a separator lower bound at least $\alpha L$ before coordinate
trimming for every genie of the stated price.  \cref{lem:ne-depth-import}
supplies a fixed $D$ and makes the permitted genie class nonempty.

Each right incidence degree is at most $B_D$, so there are at most
$B_D L$ edges.  Choose a constant degree threshold $M_0\ge8B_D/\alpha$.
Delete every coordinate whose left degree exceeds $M_0$; their number is
at most $\alpha L/8$.  Also delete coordinates with an own-port edge.
The one-stage transpose-matching estimate of \cref{cor:nb-self-demands}
bounds these by $O(L/\log d)$, uniformly over the genies in question;
eventually this is at most $\alpha L/8$.  This use of the matching estimate
does not require transpose correctness or an initial left-degree bound.
The remaining graph has maximum degree at most $M=\max(M_0,B_D)$ and
$N=\Theta(L)$.  The separator loss is at most twice the number of deleted
coordinates, giving $\alpha L/2$.  For any fixed $K$, eventually $K\le m$,
so \cref{prop:ne-separator} applies.  The code conclusion is
\cref{thm:ne-full-cone-half}, whose quantifiers already include every
induced coordinate set.
\end{informalproof}

\begin{remark}\label{rem:ne-scope}
The conclusion rules out a uniform small-component decomposition for this
family at arbitrarily small deletion cost.  It does not exhibit a high-rate
code.  In fact its simultaneous half-rate bound shows why a directed
coding theorem can remain possible after a separator approach fails.
Neither the lane lemma nor its coloring is claimed for arbitrary native
schedules.  The next construction changes the paid physical layout and
destroys almost-everywhere separation of paired ports by component colors.
\end{remark}
